\chapter{Conclusion} \label{Conclusion}

Gender is an extremely complex phenomenon in language and translation because of its interplay with both linguistic structures, e.g. word inflection and agreement, and the extralinguistic reality, i.e. the gender identity of a referent as well as societal norms and assumptions. Since linguistic structures as well as views on gender are language-dependent, translating gender is not a trivial undertaking -- for humans or machines. While issues of translating gender have been investigated both within TS and in MT, little research has tackled them beyond a binary conception of gender.

The present work thus focused on GFL as a means to refer to and represent non-binary individuals in grammatical gender languages such as German, where only the masculine and feminine gender are used in reference to people. Due to the linguistic and cultural influence of English, translators are increasingly confronted with texts depicting non-binary referents (e.g. TV series such as \textit{One Day at a Time}). While gender neutrality may be quite easily reached in English with the pronoun \textit{they} used for referents in the singular and gender-neutral alternatives to common nouns, this is not the case in German. Different strategies have been proposed which can be subsumed under the term gender-fair and include gender-inclusive and gender-neutral approaches.

This wealth of linguistic strategies is challenging both in human and MT since there is no commonly acknowledged standard for GFL, and strategy selection depends on a variety of factors. For this reason, the present work investigated the translation and PE of GFL beyond the binary in a study that was both process- and product-oriented. Twelve professional translators were recruited and asked to translate or to post-edit three English-language texts into German. The source texts were news reports about non-binary actors joining the cast of various TV series. For each text, a different GFL approach, i.e. (i) gender-neutral rewording, (ii) gender-inclusive characters, and (iii) neosystems, was to be used. To prepare for the study, participants received a handout on GFL, which could also be used during the translation and PE process.

Participants joined a video conference, which was open in the background of their computers and their screen was recorded. They were asked to complete their assignments and then they were interviewed. Cognitive processes were analysed by focusing on translation and PE times, screen activity tracking data, and the personal experiences of the study participants. Such data were elicited by means of non-participant observation, screen recordings, and retrospective interviews. Product data were analysed by annotating target texts for both the type of GFL strategy used as well as the success of its integration in the translation and PE process.

Neosystems were found to be the most difficult approach to GFL. While differences in translation and PE times were not statistically significant, screen activity tracking data (e.g. the number of searches and changes) and mistakes considerably increased when neosystems were used. Neosystems are still very rare and the study participants felt like they had to learn a foreign language, often using the provided handout to be able to use neosystems.

Participants expressed a clear preference for gender-inclusive characters, especially the gender star, combined with gender-neutral rewording. Gender-inclusive characters are a quite straightforward approach to reach gender fairness in language and their use is quite common. However, they might impact readability if frequently used within the same text. In a similar case, gender-neutral rewording is deemed a good alternative. Neutral formulations should be used carefully in reference to a single person because the target texts might sound quite unnatural or impersonal. 

The target text analysis shows that, amongst other things, gender-inclusive characters are applied differently, with language professionals wishing for clearer rules and/or standards for their use. Participants are very creative in gender-neutral rewording, often reformulating whole phrases to avoid gender. When using neosystems, participants often made mistakes, suggesting higher cognitive load and the need for further training to correctly apply them. 

PE was significantly faster than translation in the first two assignments. When using neosystems, the speed difference between tasks was not statistically significant. This seems to suggest that the higher cognitive load experienced because of the use of neosystems reduced the advantages in speed brought about by the automated translation. Participants who received the PE task also made fewer mistakes. It was also found that some gender-specific terms are more challenging than others and therefore the application of a strategy is not always straightforward or possible.

The present study has important implications for both TS and professionals. A single source text might be translated differently, depending on a large variety of factors. Such factors may include the client's specifications, target audience and culture as well as personal and societal beliefs. This also influences the selection of a specific GFL strategy, since the study shows how different strategies carry different implications. The result is that a one-size-fits-all approach to the representation of non-binary individuals is neither possible nor desirable.

GFL is constantly evolving and there are neither norms nor consensus on which and how a specific strategy is to be used; constant training and contact with the queer and non-binary community are needed. This also applies to the technology design process in which relevant stakeholders, e.g. translation professionals and non-binary people, need to be involved. Finally, current MT paradigms need to shift to allow one-to-many translations, personalisation, and strategy selection. Theories and perspectives from TS study might therefore inform the strategy selection as translation scholars and professionals are best equipped to understand the connotations of different GFL strategies in cross-linguistic scenarios.

Finally, non-binary people are still a marginalised group exposed to daily discrimination. The translation of GFL is therefore a highly sensitive area, where translators bear the responsibility for a respectful representation of marginalised groups. However, since GFL is still controversial and politically/ideologically charged, translation is best conceptualised as a process of cross-linguistic, cultural, political, and ideological negotiation. 
 
The present study is a first step towards tackling issues of non-binary representation in languages with grammatical gender. Future studies should involve a higher number of participants and also apply more precise and objective methodologies such as keylogging and eye tracking. A greater involvement of non-binary people in the research process would also be beneficial. Furthermore, future studies should investigate GFL in other language pairs, also without English as the source language and culture. Psycholinguistic studies should also shed light on the readability and comprehensibility of GFL for different groups. Finally, further exchanges on the topic between TS and MT might be beneficial to both fields, promoting multidisciplinary collaboration which should also extend to other disciplines such as sociolinguistics, human-computer interaction, philosophy (e.g. linguistic injustice), ethics and law (e.g. legal definition of discrimination and bias). 
